\scanpage{8}
% Reused typescript, reading orientation rotated 180 degrees.
\mspage{7}
\begin{center}-2.33-\end{center}
then:
\[
\begin{aligned}
\nu^{M\otimes f^*L^n}(x,\delta)&=\nu^{M\otimes f^*L^{n_0}}(x,\delta)+\nu^{f^*L^{n-n_0}}(x,\delta)\\
&=\nu^{M\otimes f^*L^{n_0}}(x,\delta)+(n-n_0)\nu^L(f(x),\delta).
\end{aligned}
\]
First of all, apply the first part of Proposition 2.17 to $\nu^{M\otimes f^*L^{n_0}}$.
We conclude that, for some $K$:
\[
\nu^{M\otimes f^*L^{n_0}}(x,\delta)\geq-K.
\]
Secondly, apply the second part of Proposition 2.17 to $\nu^L(y,\delta)$,
for $y\in S=Y^s_{(0)}(L)_k$. By\unsure{,,}Theorem 2.1, all these numbers are
strictly positive. Therefore, by the Proposition:
\[
\nu^L(y,\delta)\geq\epsilon>0
\]
for all $y\in Y^s_{(0)}(L)_k$, and $\delta\in\Delta(G)$. Therefore, if $n\geq\left[n_0+\frac K\epsilon+1\right]$,
we conclude
\[
\nu^{M\otimes f^*L^n}(x,\delta)>0
\]
whenever $x\in f^{-1}\{Y^s_{(0)}(L)\}_k$. Therefore, by Theorem 2.1, all such
$x$ are in $X^{\unsure{s}}_{(0)}(M\otimes f^*L^n)$. \uline{QED}
